\documentclass[11pt,reqno]{amsart}
\usepackage{amsmath,amssymb,amsthm}
\usepackage[margin=1.15in]{geometry}
\usepackage{booktabs}
\usepackage{microtype}
\usepackage[hidelinks]{hyperref}

\theoremstyle{plain}
\newtheorem{theorem}{Theorem}[section]
\newtheorem{lemma}[theorem]{Lemma}
\newtheorem{corollary}[theorem]{Corollary}
\theoremstyle{definition}
\newtheorem{definition}[theorem]{Definition}
\theoremstyle{remark}
\newtheorem{proposition}[theorem]{Proposition}
\newtheorem{remark}[theorem]{Remark}

\newcommand{\GF}{\mathbb{F}}
\newcommand{\rk}{\operatorname{rank}}
\newcommand{\spn}{\operatorname{span}}
\newcommand{\cost}{\operatorname{cost}}
\newcommand{\Stab}{\operatorname{Stab}}
\newcommand{\nnz}{\operatorname{nnz}}
\newcommand{\wt}{\operatorname{wt}}
\newcommand{\col}{\operatorname{col}}

\title[Exact bilinear complexity over finite fields]
      {Exact bilinear complexity over finite fields:\\
       guarantees for a descent heuristic, and a pool-free solver route}
\author{Mohamed Ali Tewfik Hamlil}
\date{\today}


\begin{document}

\begin{abstract}
The rank of a bilinear map is the number of multiplications required to compute
it, and deciding it over a finite field is NP-complete. We report on an
implementation that attacks the problem from two directions and prove what each
one guarantees. For a greedy descent we establish exactness of its first step by
Rado--Edmonds, an invariant giving soundness, a termination bound, and a pruning
lemma from which the descent's fixed point is shown to be locally optimal against
the \emph{entire} candidate pool rather than against the subset it scanned. We
prove that the improvement predicate is invariant under any subgroup of the
stabiliser of the current span, which is what makes symmetry reduction sound, and
observe that the implementation quotients by a subgroup rather than by the full
stabiliser, so that enlarging it is a question of efficiency and never of
correctness. Separately we note that the satisfiability route requires no
enumeration of rank-one maps at any point, so its space is polynomial in the
input where the combinatorial route is exponential: on $\langle 4,4,4\rangle$
the pool is $8.2$~TiB and refused, while the formula is $11$~MB and written in
$1.5$~s. A second strand minimises the additions rather than the
multiplications, and there the same Rado--Edmonds argument returns the minimum
over every change of basis rather than a good answer: on the three operators of
a published rank-23 $\langle 3,3,3\rangle$ scheme it takes $221$ nonzeros to
$128$ against $167$ for the reference implementation, and we separate the three
incomparable cost models in which such a count can be read, since on one
operator in nine the fewer nonzeros cost additions. We also record three
negative results, since each cost real time to establish.
\end{abstract}

\maketitle

\section{Introduction}

Let $\GF$ be a field and let $T$ be a bilinear map, presented as a list of
slices $T_1,\dots,T_k$, each an $n\times m$ matrix over $\GF$. The
\emph{rank} of $T$, equivalently its \emph{bilinear complexity}, is the least
$r$ for which
\[
  t_{ijl} \;=\; \sum_{s<r} a^{(s)}_i\, b^{(s)}_j\, c^{(s)}_l ,
\]
that is, the least number of multiplications sufficient to compute $T$.
Strassen's seven products in place of eight for $2\times 2$ matrix
multiplication is the origin of the subject, and determining rank in general is
open.

Two facts frame everything below. First, over a finite field the decision
problem ``is $\rk(T)\le r$?'' is NP-\emph{complete}, not merely NP-hard; the
sharper statement is due to Schaefer and \v{S}tefankovi\v{c}, who show that
tensor rank over $\GF$ is polynomial-time equivalent to the existential theory
of $\GF$. Membership in NP is not a technicality here: it is what makes a
decomposition a certificate of polynomial size, and therefore what makes the
problem expressible to a solver at all. Over $\mathbb{R}$ or $\mathbb{Q}$ no such
certificate exists, and over $\mathbb{Q}$ the problem is not known to be
decidable.

Second, the two natural algorithmic formulations differ in a way that is
invisible in their statements and decisive in their space usage. Section
\ref{sec:space} makes this precise.

Our contributions are: proofs for a descent heuristic that previously carried
measurements and one informal remark (Section~\ref{sec:descent}); the observation
that the solver formulation is pool-free and therefore polynomial in space
(Section~\ref{sec:space}); a soundness result for symmetry reduction together
with the corollary that closes the gap between it and the implementation
(Section~\ref{sec:symmetry}); a proof that the operator sparsification of the
second strand attains the minimum over every change of basis, with the ceiling
that makes its scan terminate and the cost models that keep its counts from
being read against each other (Section~\ref{sec:sparsify}); and three negative
results (Section~\ref{sec:neg}).

\section{Preliminaries}

Write $\cost(X)=\sum_i \rk(X_i)$ for a list $X$ of matrices. A rank-one
bilinear form is exactly one multiplication, so $\cost$ counts multiplications.

\begin{definition}
A list $S$ is \emph{feasible} for $T$ if $\spn(S)\supseteq\spn(T)$.
\end{definition}

Feasibility, rather than equality of spans, is the right notion: $S$ need only
\emph{generate} $T$, so a larger spanning set of smaller total rank is a better
answer. Karatsuba's five products for a four-coefficient product is exactly such
a case, and it is why $|S|$ may exceed $|T|$.

We use throughout that $\spn(T)\setminus\{0\}$ under linear independence is a
vector matroid, and that $w(M)=\rk(M)$ is a non-negative integer weight on it.

\section{A descent, and what it guarantees}\label{sec:descent}

The method has three steps. Step 1 chooses a basis of $\spn(T)$ of least total
rank. Steps 2 and 3 relax feasibility from ``basis'' to ``spanning set'' and
descend over a pool $G$ of candidate rank-one maps, adopting any candidate whose
addition lowers $\cost$. The two later steps are the same procedure differing
only in $G$: step 2 uses the rank-one terms of $T$ itself, step 3 every rank-one
map of the shape, one per scalar class, of which there are
$(p^{n}-1)(p^{m}-1)/(p-1)^2$.

\begin{theorem}[Step 1 is exact; \cite{nakatsukasa2017}, Theorem 2.1]\label{thm:greedy}
Sorting $\spn(T)\setminus\{0\}$ by ascending rank and keeping each element that
preserves independence returns a basis of $\spn(T)$ of minimum total rank.
\end{theorem}

\begin{proof}
The procedure is the matroid greedy. By Rado--Edmonds it returns a
minimum-weight basis of any matroid under any non-negative weight function. The
matroid is the vector matroid of $\spn(T)$ and the weight is $\rk$.
\end{proof}

\begin{remark}[Attribution, and a stronger statement]
Theorem~\ref{thm:greedy} is not new and is stated here only because everything
below depends on it. Nakatsukasa, Soma and Uschmajew pose exactly this
minimisation, give exactly this greedy as their Algorithm~1, and make the matroid
observation in the same words, calling it a greedy algorithm for an infinite
matroid of finite rank \cite{nakatsukasa2017}. Their Corollary~1 is stronger than
the remark below: not only is the optimal \emph{value} independent of tie-breaks,
but any other basis of lowest total rank takes the same ranks up to permutation,
so the whole multiset is an invariant of $\spn(T)$. They also record the caveat
that the greedy's oracle, finding a lowest-rank element outside the span so far,
is itself \textsc{NP}-hard, so exactness of the \emph{choice} is not tractability
of the \emph{problem}. What is new here is the instantiation of that oracle over
$\GF(p)$ by enumeration, which their nuclear-norm route has no analogue for.
\end{remark}

The \emph{value} is therefore independent of how ties are broken, all
minimum-weight bases of a matroid having equal weight; the \emph{basis} need not
be. Fixing the enumeration order accordingly buys reproducibility, not
optimality.

\begin{theorem}[Soundness]\label{thm:sound}
Every $S$ the descent holds is feasible for $T$.
\end{theorem}

\begin{proof}
By induction. The initial $S_0$ is a basis of $\spn(T)$, so $\spn(S_0)=\spn(T)$.
A step replaces $S$ by a basis of $\spn(S\cup\{\varphi\})$, and
$\spn(S\cup\{\varphi\})\supseteq\spn(S)\supseteq\spn(T)$ by hypothesis.
\end{proof}

Consequently a result is never wrong, only possibly large.

\begin{theorem}[Termination]\label{thm:term}
The descent performs at most $\cost(S_0)-\rk(T)$ adoptions.
\end{theorem}

\begin{proof}
A candidate is adopted only when $\cost$ strictly decreases. By
Theorem~\ref{thm:sound} every intermediate $S$ is feasible, so
$\cost(S)\ge \rk(T)$. A strictly decreasing sequence of integers bounded below is
finite, of length at most the initial gap.
\end{proof}

The next lemma is what licenses the implementation to discard candidates
permanently rather than for a single pass, and it is the hypothesis of
Theorem~\ref{thm:oneopt}.

\begin{lemma}[Pruning is sound]\label{lem:prune}
If $\varphi\in\spn(S)$ then $\varphi$ cannot lower $\cost$, now or at any later
stage of the descent.
\end{lemma}

\begin{proof}
Each $S$ is a minimum-weight basis of its own span, being the output of
Theorem~\ref{thm:greedy} applied to that span. If $\varphi\in\spn(S)$ then
$\spn(S\cup\{\varphi\})=\spn(S)$, whose minimum-weight basis has weight
$\cost(S)$; there is no decrease. By Theorem~\ref{thm:sound} the span only grows,
so $\varphi\in\spn(S')$ for every later $S'$ and the argument repeats.
\end{proof}

\begin{theorem}[Local optimality over the whole pool]\label{thm:oneopt}
At the fixed point, no $\varphi\in G$ lowers $\cost$, including those never
tested.
\end{theorem}

\begin{proof}
Each $\varphi\in G$ was either tested and failed to improve, or discarded, in
which case Lemma~\ref{lem:prune} shows it could not have improved.
\end{proof}

Theorem~\ref{thm:oneopt} is strictly stronger than the loop's stopping condition,
which speaks only of the candidates that survived pruning. It is also the claim
most exposed to an implementation error, since a defect in the pruning would
produce a quietly worse answer rather than a visible failure; it is therefore
asserted directly by a test, which rescans the entire pool at the fixed point.

\begin{remark}
No approximation ratio follows. Theorem~\ref{thm:oneopt} concerns single
additions only. A move exchanging two elements at once lies outside the
neighbourhood and can pay: $\langle 2,2,2\rangle$ sits at a fixed point of cost
$8$ while its rank is $7$, which is why escaping such plateaus is a separate
method rather than a parameter of this one.
\end{remark}

\begin{remark}
Step 3 cannot do worse than step 2. Step 2's pool consists of rank-one terms of
$T$, each a scalar multiple of an element of step 3's pool, and $\varphi$ and
$c\varphi$ span the same line. The ordering of the steps therefore costs time
and not quality.
\end{remark}

\section{Symmetry reduction}\label{sec:symmetry}

Let $\sigma=(\mu,\nu)$ with $\mu,\nu$ invertible act on matrices by
$M\mapsto \mu^{\mathsf T} M\nu$. This is the rank-preserving action under which
a decomposition of $T$ is carried to another decomposition.

\begin{theorem}[Orbit invariance]\label{thm:orbit}
If $\sigma(\spn S)=\spn S$ then $\varphi$ lowers $\cost(S)$ if and only if
$\sigma(\varphi)$ does.
\end{theorem}

\begin{proof}
Invertibility of $\mu,\nu$ gives $\rk(\mu^{\mathsf T}M\nu)=\rk(M)$, so $\sigma$
preserves weights and $\cost(X)=\cost(\sigma(X))$ for every list $X$. Being a
linear isomorphism fixing $\spn(S)$, it carries $\spn(S\cup\{\varphi\})$ onto
$\spn(S\cup\{\sigma\varphi\})$. By Theorem~\ref{thm:greedy} minimum weight is a
function of the span alone, so the two additions have equal cost.
\end{proof}

Hence one representative per orbit of $\Stab(\spn S)$ suffices. Note that the
hypothesis concerns the stabiliser of the \emph{current} span and not of $T$: a
quotient computed once and reused as $S$ moves is unsound, not merely weaker.

\begin{corollary}\label{cor:subgroup}
The conclusion of Theorem~\ref{thm:orbit} holds for every subgroup
$H\le \Stab(\spn S)$, with orbits at least as fine.
\end{corollary}

\begin{proof}
The proof of Theorem~\ref{thm:orbit} uses only that each element preserves rank
and fixes $\spn(S)$, which every element of a subgroup does. Finer orbits yield
more representatives, hence a smaller reduction and never a wrong answer.
\end{proof}

Corollary~\ref{cor:subgroup} is what reconciles the theory with the
implementation. On the matrix multiplication path the ambient group is supplied
as nine generators, and the code retains those that fix the span; the quotient is
therefore by the subgroup they generate and not by the full stabiliser. The
corollary shows this is sound, and locates the risk in any future improvement:
computing a larger subgroup changes performance and cannot change an answer.

\begin{remark}
The action above is not a symmetry a satisfiability solver can exploit
internally. Every established SAT symmetry method acts on permutations of
literals commuting with negation, whereas $\sigma$ mixes coordinates linearly.
Only the symmetric group on the $r$ summands is expressible in conjunctive normal
form. Group-based reduction must therefore enter from outside the solver, as a
partition of the search into independent instances, one per orbit of the first
term.
\end{remark}

\section{Two formulations, and why their space differs}\label{sec:space}

The combinatorial formulation searches for a subspace of dimension $r$
possessing a basis of rank-one maps, choosing successive basis elements from an
explicit pool. The satisfiability formulation introduces one variable per
coordinate of each factor and constrains their products to reproduce $T$.

The two decide the same question and their space requirements are not
comparable. The pool has $(p^{n}-1)(p^{m}-1)/(p-1)^2$ elements, exponential in
the axis lengths; the formula has $r(n_1+n_2+n_3)+r\,n_1n_2+r\,n_1n_2n_3$
variables, polynomial in the size of $T$. The rank-one condition in the second is
imposed by the shape of the formula rather than by selection from a list, so no
enumeration occurs at any point.

\begin{table}[h]
\centering
\begin{tabular}{lrr}
\toprule
$\langle 4,4,4\rangle$ at $r=47$ & combinatorial & satisfiability \\
\midrule
objects held      & $4{,}294{,}836{,}225$ & --- \\
space             & $8.2$ TiB (refused)   & $11$ MB \\
time to produce   & ---                   & $1.5$ s \\
\bottomrule
\end{tabular}
\vspace{0.5em}
\caption{The same tensor and the same $r$, both instruments.}
\end{table}

The distinction is not incidental to the implementations. It follows from the
formulations: ``choose the $i$-th rank-one map'' presupposes an enumeration,
while ``these coordinates multiply out to $T$'' does not.

\begin{remark}
Polynomial space is not tractability. For the combinatorial route a lazily
generated pool would change space and not time, converting an immediate failure
into a computation that does not terminate in practice. For the solver route the
distinction is the whole point, since the search itself is delegated.
\end{remark}

\begin{remark}
Membership in FNP is realised literally. On a satisfying assignment the
implementation reconstructs the rank-one matrices, multiplies them out, and
compares against $T$, rejecting the answer if they differ. A positive answer
therefore does not rest on trusting the solver or the encoding. A negative answer
may be accompanied by a DRAT refutation checked by an independent program.
\end{remark}

\section{Two slices: a theorem that does not survive the field}\label{sec:pencil}

A tensor with two slices is a matrix pencil $A + xB$, and Kronecker's theory
classifies such a pencil up to strict equivalence by its \emph{minimal indices}
$\varepsilon_i, \eta_j$ and the \emph{elementary divisors} of its regular part.
Ja'Ja' \cite{jaja1979}, and independently Grigoriev \cite{grigoriev1978}, give
the rank from that data: over an algebraically closed field,
\begin{equation}\label{eq:jaja}
  R(T) \;=\; \sum_{\varepsilon_i > 0}(\varepsilon_i + 1)
          \;+\; \sum_{\eta_j > 0}(\eta_j + 1)
          \;+\; r \;+\; \delta,
\end{equation}
where $r$ is the size of the regular part and $\delta$ counts the elementary
divisors of degree at least two. All of the data in \eqref{eq:jaja} is
computable over $\mathrm{GF}(p)$ in polynomial time, which is what
\texttt{pencil\_rank/} does: the singular structure from the ranks of one block
system, the divisors from a Smith diagonal over $\mathrm{GF}(p)[x]$ taken twice,
forwards for the finite divisors and reversed for those at infinity.

\begin{remark}[The mistake worth recording]
Applying \eqref{eq:jaja} with the elementary divisors computed \emph{over
$\mathrm{GF}(p)$} yields neither side. It is not the rank over the closure,
where an irreducible of degree $d$ splits into $d$ distinct linear forms and
stops contributing to $\delta$; and it is not the rank over $\mathrm{GF}(p)$.
For $(I_4, C)$ over $\mathrm{GF}(2)$ with $C$ the companion matrix of $x^4 + x +
1$, that reading gives $5$. An exhaustive search proves no decomposition with
five products exists, walking $1\,897\,576$ nodes to the end of the tree, and
exhibits one with six.
\end{remark}

\begin{proposition}\label{prop:closure}
Let $T$ be a two-slice tensor over $\mathrm{GF}(q)$. Then $R_{\mathrm{GF}(q)}(T)
\ge R_{\overline{\mathrm{GF}(q)}}(T)$, and the right-hand side is given exactly
by \eqref{eq:jaja} with the elementary divisors taken over the closure. If
moreover $T$ is diagonalisable over $\mathrm{GF}(q)$, meaning it has no minimal
indices and every elementary divisor is a linear form of multiplicity one, the
inequality is an equality.
\end{proposition}

\begin{proof}
A decomposition over $\mathrm{GF}(q)$ is a decomposition over any extension, so
enlarging the field cannot raise the rank; the value over the closure is
\eqref{eq:jaja} by \cite{jaja1979}. For the second claim, diagonalisability
gives $P, Q$ over $\mathrm{GF}(q)$ with $P A Q$ and $P B Q$ both diagonal, and
the $r$ matrices $e_i e_i^{\mathsf T}$ are rank one, rational over
$\mathrm{GF}(q)$, and span a space containing both. So $R \le r$, and $r$ is the
value of \eqref{eq:jaja} here.
\end{proof}

The gap is a field-size effect rather than a defect in \eqref{eq:jaja}, and it
is not new. Sumi, Miyazaki and Sakata \cite{sumi2009} record the condition
exactly. Writing $k$ for the number of invariant polynomials of $A$ that do not
factor into \emph{distinct} linear factors over $K$, their Theorem~3.3, after
Ja'Ja', gives $\rk_K(E^n; A) \le n + k$ \emph{provided}
$\mathrm{Card}(K) \ge \deg p_1(A)$; their Theorem~3.5 gives
$\rk_K(E^n; A) \ge n + k$ with no hypothesis at all, where Ja'Ja' had the
reverse inequality only when $p_1(A)$ splits; and their Proposition~3.4 exhibits
$(E_3, A)$ over $\mathrm{GF}(2)$, with $A$ the companion matrix of $x^3 + x + 1$,
of rank at least $5$, so the cardinality hypothesis cannot be dropped.

That proposition is one of the twelve pencils settled by exhaustion here, and
this work measured it at exactly $5$ against a count of $4$ before finding the
statement. So the measurements confirm a theorem rather than exhibiting a gap,
and \texttt{pencil\_rank} computes $n + k$ and tests the cardinality condition:
where it holds the rank is settled outright, and where it fails $n + k$ remains
a proved lower bound over every field. What is left open is the cost \emph{below}
the threshold, which is where the three short rows of that table live and where
neither theorem applies.

\section{The other cost: sparsifying the operators}\label{sec:sparsify}

A decomposition fixes the multiplications and is silent about the additions,
which are set by how many nonzero entries the operators that encode and decode
it carry. An operator $U$ may be replaced by $U\Phi$ for any invertible $\Phi$,
which computes the same bilinear map in another basis and in general costs fewer
additions, so a second minimisation sits beside the first.

\begin{definition}[Operator sparsification]\label{def:sparsify}
Given $U\in\GF^{N\times d}$ of full column rank, find an invertible
$\Phi\in\GF^{d\times d}$ minimising $\nnz(U\Phi)$.
\end{definition}

The columns of $U\Phi$ are a basis of $\col(U)$, every basis of $\col(U)$ is the
column list of $U\Phi$ for exactly one such $\Phi$, and $\nnz$ of a matrix is the
sum of the Hamming weights of its columns. Definition~\ref{def:sparsify} is
therefore the request for a basis of $\col(U)$ of least total weight: the problem
of Theorem~\ref{thm:greedy} with $\wt$ where that one has $\rk$, and its solution
is the same theorem. The hardness surrounding it is a different literature
entirely, and is stated first, because a section proving a minimum without
saying what is hard invites the reading that the problem was easy.

\emph{Four names, and two unrelated hardness proofs.} The inner question, the
lightest nonzero vector of a subspace, is one object the literature names four
ways: the \emph{minimum distance} of a linear code over $\GF(q)$; the
\emph{sparsest vector in a subspace} over $\mathbb{R}$ or $\mathbb{Q}$
\cite{qu2020}; the \emph{spark} of a matrix \cite{tillmann2014}; and
the \emph{girth}, the shortest circuit, of a represented matroid. The
identification is stated outright in the dual formulation of
\cite{bhattacharyya2021}: the minimum distance of the code generated by $A$ is
the size of the smallest linearly dependent set among the columns of $A^{\perp}$,
which is the shortest circuit of the binary matroid they represent. The greedy's
first call below is exactly this problem and its later calls are the relative
version of it, so everything here applies to both.

Over $\mathbb{Q}$ the sparsest-basis problem is NP-hard \cite{mccormick1983}, as
is the sparse null space basis problem, and that one even when the optimal value
is supplied \cite{colemanpothen1986}. It admits no approximation within
$2^{\log^{1/2-o(1)} n}$ \cite{gottlieb2010}, which is the result the oracles of
\cite{beniamini2020} are built around, and deciding whether a rational matrix has
a circuit of size at most $k$ is \emph{strongly} NP-complete \cite{tillmann2014}.
The sharpest statement over $\mathbb{R}$ is that the sparsest vector cannot be
approximated within any constant factor \cite{sparsevectorfocs2025}.

Over a finite field the minimum distance problem is NP-complete over $\GF(2)$
\cite{vardy1997}; over \emph{every} finite field it is hard to approximate within
any constant under randomised reductions, and within
$2^{\log^{1-\varepsilon} n}$ under quasi-polynomial ones \cite{dumer2003}; and
parameterised by the distance bound it is
$W[1]$-hard to approximate within any constant, over $\GF(2)$
\cite{bhattacharyya2021} and over any fixed finite field \cite{bennett2023}.

Neither version is known to reduce to the other, in either direction. The two
proofs are independent, the inapproximability results are twenty-two years apart,
and the obstruction has a name: over a finite field one step is free from the
second generalised Hamming weight bound, and over $\mathbb{R}$ the same step took
a chaining argument and Littlewood--Offord anticoncentration.

\begin{remark}[Two attributions that are routinely got wrong]
\cite{berlekamp1978} proves \emph{coset weight} and \emph{subspace weight}
NP-complete, the second being ``is there a codeword of weight exactly $w$''; its
own abstract says the result strongly suggests but does not rigorously imply the
minimum distance case, which \cite{vardy1997} closed nineteen years later. And
\cite{downey1999} proves maximum-likelihood decoding and weight distribution
$W[1]$-hard, and not minimum distance, which it leaves open. Both are commonly
given as citations for minimum distance, and were so given here first.
\end{remark}

\begin{remark}[Why the exact method below contradicts none of it]\label{rem:notcontra}
Over $\GF(q)$ the implementation lists the $q^{d}$ elements of the space and
takes them in ascending weight. That is exponential in $d$ and not polynomial in
the input, so it sits where the hardness does not bite: the finite-field
operators here have $d$ equal to $5$ or $6$, and $q^{d}$ is $32$ or $64$. Nor is
the factor an artefact of a naive implementation. Under the strong exponential
time hypothesis there is no $q^{(1-\varepsilon)d}$ algorithm for the minimum
distance problem over any finite field, for a code of $q^{d}$ codewords
\cite{stephensdavidowitz2019}, which is a lower bound of exactly this walk's
shape. Over $\mathbb{Q}$ the space cannot be listed at all, which is what forces
the scan of Theorem~\ref{thm:sparseceiling}.
\end{remark}

\emph{The exact oracle, and why the greedy is minimal.}

\begin{theorem}[Sparsification by the matroid greedy is exact]\label{thm:sparseexact}
Let $\mathcal{C}=\col(U)$ have dimension $d$. Scanning $\mathcal{C}\setminus\{0\}$
by ascending Hamming weight and keeping each vector that preserves independence
returns a basis of $\mathcal{C}$ of minimum total weight, and the $\Phi$ carrying
the columns of $U$ to that basis attains the minimum of
Definition~\ref{def:sparsify}.
\end{theorem}

\begin{proof}
The two minimisations are one, the columns of $U\Phi$ ranging over exactly the
bases of $\mathcal{C}$ as $\Phi$ ranges over the invertible matrices. Linear
independence in $\mathcal{C}$ is the vector matroid of any matrix whose columns
span it \cite[Prop.~1.1.1]{oxley}, and $\wt$ is a non-negative integer weight on
that matroid. The procedure is the matroid greedy, which by Rado--Edmonds returns
a minimum-weight basis of any matroid under any such weight
\cite[Lem.~1.8.3]{oxley}. This is Theorem~\ref{thm:greedy} again, over a
different matroid and under a different weight.
\end{proof}

\begin{remark}[Attribution]
The composition is not new. The driver is \cite{gottlieb2010}'s, reproduced by
\cite{beniamini2020} as its Algorithm~2, and the oracle it calls is their
Problem~2.15, the lightest vector of the space lying outside the span of what is
already settled; given such a subroutine, they state, Algorithm~2 returns an
exact solution. What is supplied here is that oracle, over $\GF(q)$ by the
enumeration of Remark~\ref{rem:notcontra} and over $\mathbb{Q}$ by the scan
below. Their own two exact oracles, Algorithms~3 and~4, were implemented here as
well; they return the same counts on every operator measured, between $88$ and
$343$ times more slowly, so what separates the methods is cost and never the
answer. It is worth recording that the same oracle is published a second time
under another name: the relative generalised Hamming weight of \cite{sanjose2025}
at $r=1$ is Problem~2.15 exactly, with a Brouwer--Zimmermann pruning already on
it, and neither literature cites the other.
\end{remark}

Over $\mathbb{Q}$ the space is infinite and the enumeration is unavailable. A
vector of weight $w$ is supported on $w$ of the $N$ coordinates, so ascending
support size is ascending weight, and restricted to a support $T$ the vectors of
$\mathcal{C}$ solve one small homogeneous system: in reduced row echelon form a
vector's entry on a pivot coordinate is that echelon row's coefficient, so the
unknowns are the echelon rows whose pivot lies inside $T$, never more than $|T|$
of them, and the equations are the coordinates that are neither pivots nor in
$T$. Walking the supports in ascending size therefore instantiates the greedy.
What is not obvious is that the walk ends.

\begin{theorem}[The greedy never exceeds the Singleton ceiling]\label{thm:sparseceiling}
Every vector the greedy of Theorem~\ref{thm:sparseexact} takes has weight at most
$N-d+1$.
\end{theorem}

\begin{proof}
Fix an information set, that is $d$ coordinates on which $\mathcal{C}$ projects
isomorphically. The systematic basis it determines consists of $d$ vectors, each
carrying a $1$ on its own coordinate of the set, a $0$ on the other $d-1$, and
arbitrary entries on the remaining $N-d$; each therefore has weight at most
$N-d+1$. At any point of the greedy fewer than $d$ vectors are held, so their
span cannot contain all $d$ systematic vectors and at least one lies outside it.
The greedy takes a lightest vector outside the current span, whose weight is
consequently at most $N-d+1$.
\end{proof}

\begin{remark}
Theorem~\ref{thm:sparseceiling} is the Singleton bound applied at every step of
the greedy rather than only at the first, and it settles a worry rather than
improving a bound. Reading an oracle as ``search the column subsets of one fixed
size'' invites the fear that once some vectors are settled the lightest vector
outside their span is heavier than the scan has reached, and is missed. It cannot
be. On the three operators of Table~\ref{tab:plinopt}, where $N=23$ and $d=9$,
the ceiling is $15$ and the largest weights actually taken were $6$, $6$ and $7$.
\end{remark}

The scan costs $O\!\left(\sum_{w\le W}\binom{N}{w}\,(w^{3}+wN)\right)$ time, with
$W$ the largest weight the greedy takes and $W\le N-d+1$ by
Theorem~\ref{thm:sparseceiling}, and $\Theta(dN)$ space, the supports being
walked rather than listed.

\begin{remark}[Where the scan stops, and what answers past it]
Theorem~\ref{thm:sparseceiling} is an upper bound and the scan carries no lower
bound, so it halts on having collected a basis rather than on a proof that
nothing lighter exists, and its cost is set by the ceiling and not by the answer.
On a $16$-dimensional subspace of $\mathbb{Q}^{49}$, the left operator of a
published rank-49 $\langle 4,4,4\rangle$ scheme, seven of the sixteen vectors
need weight $6$ or more and the widest level the walk could reach prices at
$1.4$ PiB, so the implementation prices it in advance and refuses in
milliseconds rather than running. Brouwer--Zimmermann
\cite{zimmermann1996,lisonek2016}, coding theory's standard method for the same
subproblem, holds a lower bound from several disjoint information sets alongside
the upper one and stops early; it is the first thing to try. Meanwhile a linear
programme answers that operator in $0.34$ s at $100$ nonzeros, which is an upper
bound and not the minimum: it would be the minimum if the operator's column
matroid were regular, and it is not, a sixteen-column minor having determinant
$-2$.
\end{remark}

\emph{Against the reference implementation.} PLinOpt \cite{plinopt} is the
reference implementation in this problem area. It reaches sparsity by a
different route, sparse QLUP elimination and a coefficient search bounded to
four rows of support with a set of eleven coefficients, and it is a heuristic
where Theorem~\ref{thm:sparseexact} is a minimum.
Table~\ref{tab:plinopt} puts the two on the three operators of a published
rank-23 $\langle 3,3,3\rangle$ scheme.

\begin{table}[ht]
\centering
\begin{tabular}{llrrrr}
\toprule
operator & shape & as given & \cite{plinopt} & minimum & seconds \\
\midrule
$L$   & $23\times 9$ & $69$  & $59$  & $\mathbf{43}$  & $0.338$ \\
$R$   & $23\times 9$ & $66$  & $53$  & $\mathbf{42}$  & $0.327$ \\
$P$   & $9\times 23$ & $86$  & $55$  & $\mathbf{43}$  & $0.434$ \\
\midrule
total &              & $221$ & $167$ & $\mathbf{128}$ & \\
\bottomrule
\end{tabular}
\vspace{0.5em}
\caption{Nonzeros in the three operators of \texttt{3x3x3\_23\_Grey-221}. The
minimum column is the minimum over every invertible $\Phi$, and it is the
minimum because Theorem~\ref{thm:sparseexact} says so and not because nothing
better was found.}\label{tab:plinopt}
\end{table}

Four things qualify that table, and each is the kind of qualification a count
invites.

\emph{The two tools are not doing the same job.} PLinOpt is downstream of a
decomposition: it is handed operators and produces a straight-line program, with
in-place accumulation, Tellegen transposition and compaction around it. The work
here searches for the decomposition and sparsifies as a second strand, and finds
no common subexpressions at all. The table compares one column of that pipeline
and not the pipeline.

\emph{The middle column is carried and not re-derived here.} Those three numbers
were measured once, on one machine, with the reference sparsifier built from
upstream and given its best settings first: raising its coefficient bound from
the default of eleven to twenty and then to forty leaves it where the default
does. That binary is not built in the checkout vendored alongside this work, so
re-running the column means building the tool first. A carried number reported as
a fresh one is the failure this sentence exists to prevent.

\emph{$P$ has to be handed over transposed.} It is $9\times 23$ and computes $9$
outputs from $23$ products, so the admissible change of basis acts on its rows
and the object is its row space. Given the stored orientation the reference
sparsifier reports $9$ nonzeros, having recombined the $23$ products, which no
change of basis on the outputs may do; its own output flags the alternative
reading. Transposed it gives the $55$ in the table. The mistake would have
produced a false defeat rather than a false victory, which is the direction an
author is least likely to check.

\emph{Sparsity is trivial to improve by returning a different matrix}, so each
answer was re-checked outside the tool in exact rational arithmetic: rank $9$
before, rank $9$ after, and rank $9$ with the two stacked, so the space returned
is the space that was given.

One row of the comparison goes the other way and is the one that matters most
here. Over $\GF(2)$ and $\GF(3)$ the reference sparsifier returns the zero matrix
and exits successfully, so on the finite fields where every operator the descent
of Section~\ref{sec:descent} emits lives, there is no column to compare against.

\emph{Fewer nonzeros is not fewer additions.} Nonzeros are a proxy for additions
and the proxy is imperfect, so a count means nothing until the model it is
counted in is named. There are three, and they are not comparable.

\begin{table}[ht]
\centering
\begin{tabular}{clc}
\toprule
 & how the additions are counted & basis change \\
\midrule
(a) & $\nnz(U)-\mathrm{rows}(U)$ and so on, no reuse of intermediates & no \\
(b) & the same formula on a sparsified triple, basis maps charged apart & yes \\
(c) & the length of a straight-line program, which may fall below (a) & no \\
\bottomrule
\end{tabular}
\vspace{0.5em}
\caption{Strassen's $18$ additions and Winograd's $15$ are model (a); the
alternative-basis line \cite{karstadt2017,beniamini2020,holtz2025} is (b), where
the basis maps cost $O(n^{2}\log n)$ per recursion level and so never reach the
leading coefficient; \cite{plinopt} and the $\langle 3,3,3\rangle$ record chain
are (c).}\label{tab:models}
\end{table}

This work measures $\nnz$, so it lives in (a) and (b) and structurally cannot see
(c). For $\langle 2,2,2\rangle$ that is enough to close the question: fifteen
additions are necessary in the standard basis for every rank-$7$ algorithm, over
an arbitrary ring \cite{probert1976,bshouty1995}; changing basis is what gets
past that bound, and twelve is optimal once it does
\cite[Thm.~1.2]{karstadt2017}. The fixtures here reach twelve, $3+3+6$ against
the published $5+5+8$, which takes Strassen's leading coefficient from $7$ to
$5$. No search over that fixture can win, and anything that appears to has
mismeasured.

Crossing (b) with (c) is what the exact stage is for, its output being available
to something that does find common subexpressions. Table~\ref{tab:cse} does that
on three published rank-23 $\langle 3,3,3\rangle$ schemes.

\begin{table}[ht]
\centering
\begin{tabular}{lrrrrr}
\toprule
 & \multicolumn{2}{c}{as published} & \multicolumn{2}{c}{after the exact stage} & \\
\cmidrule(lr){2-3}\cmidrule(lr){4-5}
scheme & $\nnz$ & program & $\nnz$ & program & \\
\midrule
\texttt{Grey-221}      & $221$ & $81$ & $128$ & $62$ & $-23\%$ \\
\texttt{3x3x3\_23\_58} & $175$ & $57$ & $119$ & $56$ & $-2\%$  \\
\texttt{3x3x3\_23\_55} & $177$ & $56$ & $117$ & $52$ & $-7\%$  \\
\bottomrule
\end{tabular}
\vspace{0.5em}
\caption{Model (b) crossed with model (c): the nonzero count, and the length of
the straight-line program a subexpression pass finds from it. The last column is
the change in that length.}\label{tab:cse}
\end{table}

The margin collapses on the two schemes somebody has already worked over, and on
one operator of the nine it reverses: \texttt{3x3x3\_23\_58\_L} arrives with $49$
nonzeros, the exact stage takes it to $43$, and the subexpression pass then finds
$18$ additions where it had found $16$ on the original. Six more zeros cost two
additions. A basis of least weight is not a basis with the most shared
subexpressions, and nothing in Theorem~\ref{thm:sparseexact} says otherwise.

\begin{remark}[What the right-hand columns are not]
Three cautions, without which those figures are not numbers. The subexpression
pass is a randomised search taking no seed: four consecutive runs at its default
effort, on one unchanged file, gave $22$, $22$, $22$, $21$. Every figure in
Table~\ref{tab:cse} is at a loop count of $10\,000$, five runs a side, best of
each, the same effort on both sides, since giving this side more would be the
easiest way to fake the comparison. Next, a minimum-weight basis is not unique:
the count is the minimum and that is proved, but which basis of that weight comes
back depends on the order the scan walks its supports, and different ones give
the downstream pass different work. An earlier run returning a different
$43$-nonzero basis of \texttt{Grey-221}'s $L$ reached $17$ where this one reaches
$18$. A figure in those columns is a figure about one basis and not about the
method, which promises the $\nnz$ column only. Last, two of the nine operators
have identical row-weight profiles and identical numbers at every stage; they are
almost certainly one operator relabelled, so the nine are eight data points.
\end{remark}

\begin{remark}[What the $62$ may not be read against]
The right-hand column of Table~\ref{tab:cse} is (b) crossed with (c), which is a
fourth thing again. Three changes of basis, one per operator, go uncharged in it:
they are what the sparsification produced, and the alternative-basis accounting
charges them separately at $O(n^{2}\log n)$ per level. The record for
$\langle 3,3,3\rangle$ at rank 23 in model (c) is $55$ additions
\cite{karunaratne2026}, in the standard basis, with nothing to charge. On that
same scheme the table says $52$, and the arithmetic a reader will do is not
available: the $52$ leaves three $9\times 9$ basis changes uncharged, the $55$ has
none, and nothing here computes the leading coefficient of an alternative-basis
algorithm with common subexpressions, which is the quantity through which the two
would have to meet. Nor do the two records for that shape compete: $61$ linear
operations in model (b) \cite{beniamini2020}, whose table \cite{holtz2025}
reproduces, and the $55$ in model (c). The $55$ implies a leading coefficient
near $4.93$, below \cite{holtz2025}'s lower bound of $5$, which is not a
contradiction, that bound being stated in model (b) and a program with common
subexpressions not being bound by $\nnz-\mathrm{rows}$. What the $56$ in the same
row of Table~\ref{tab:cse} does say is that the subexpression pass at that effort
roughly reproduces its own shipped record, one addition above it, which is the
check that the left-hand columns are measured properly at all.
\end{remark}

\section{Negative results}\label{sec:neg}

Each of the following cost time to establish and is recorded so that it is not
established again.

\emph{Fixed-$k$ search with commitment.} A finder that asks only questions it
expects to be satisfiable was built on an assumed asymmetry between the cost of
acceptance and refutation, quoted at two orders of magnitude. Measured with
matched flags, the asymmetry is approximately unity; the original figures
compared instances encoded with and without symmetry breaking. The method is
dominated by the descent on every fixture measured.

\emph{Deeper cube splitting.} Partitioning a refutation by orbits of the first
two terms rather than the first was measured at $2989.8$ s against $108.2$ s for
the undivided instance, with the longest single part exceeding the whole. Even
unbounded parallelism does not recover this.

\emph{The expensive step of the descent.} Step 3 improved the answer on two of
four fixtures, by one product each time, at between $58$ and $184$ times the cost
of the first two steps combined. Effort spent accelerating it is spent on the
part that mostly does not pay.

\section{Limitations}

The descent carries no approximation guarantee, and \cite{swernofsky2018} is
why it is unlikely to acquire a good one: approximating the rank of a
three-tensor within $1 + 1/1852$ is \textsc{NP}-hard \emph{over any field}, so
the obstruction is not an artefact of working over $\GF(p)$ and it bounds every
method here at once. The exhaustive route is complete but its cost is
$\binom{|G|}{\,r-\dim T\,}$, which is prohibitive beyond small shapes. The
sparsification of Section~\ref{sec:sparsify} is exact but priced by
Theorem~\ref{thm:sparseceiling}'s ceiling rather than by its answer, so it does
not finish on an operator whose greedy weights are large, and what it minimises
is a proxy that the same section measures going the wrong way on
one operator in nine. Reported timings are machine-dependent and are not
reproducible in continuous integration; counts are, and are asserted there.
Nothing here settles the bilinear rank problem.

\begin{thebibliography}{99}
\bibitem{schaefer} M.~Schaefer and D.~\v{S}tefankovi\v{c},
  \emph{The complexity of tensor rank}, Theory of Computing Systems, 2018.
\bibitem{hastad} J.~H\r{a}stad, \emph{Tensor rank is NP-complete},
  Journal of Algorithms 11 (1990), 644--654.
\bibitem{bdez} R.~Barbulescu, J.~Detrey, N.~Estibals and P.~Zimmermann,
  \emph{Finding optimal formulae for bilinear maps}, WAIFI 2012.
\bibitem{jaja1979} J.~Ja'Ja', \emph{Optimal evaluation of pairs of bilinear
  forms}, SIAM Journal on Computing 8 (1979), 443--462.
\bibitem{grigoriev1978} D.~Yu.~Grigoriev, \emph{Some new bounds on tensor rank},
  LOMI preprint, 1978.
\bibitem{sumi2009} T.~Sumi, M.~Miyazaki, T.~Sakata, \emph{Rank of 3-tensors
  with 2 slices and Kronecker canonical forms}, Linear Algebra and its
  Applications \textbf{431} (2009), 1858--1868.
\bibitem{nakatsukasa2017} Y.~Nakatsukasa, T.~Soma, A.~Uschmajew,
  \emph{Finding a low-rank basis in a matrix subspace}, Mathematical
  Programming \textbf{162} (2017), 325--361.
\bibitem{swernofsky2018} J.~Swernofsky, \emph{Tensor rank is hard to
  approximate}, APPROX/RANDOM 2018, LIPIcs \textbf{116}, 26:1--26:9.
  (Factor $1 + 1/1852$, over any field.)
\bibitem{oxley} J.~Oxley, \emph{Matroid Theory}, 2nd ed., Oxford, 2011.
  (Rado--Edmonds, the matroid greedy.)
\bibitem{gottlieb2010} L.-A.~Gottlieb and T.~Neylon, \emph{Matrix sparsification
  and the sparse null space problem}, APPROX/RANDOM 2010.
  (The greedy driver the sparsest-vector oracles are oracles for.)
\bibitem{beniamini2020} G.~Beniamini, N.~Cheng, O.~Holtz, E.~Karstadt and
  O.~Schwartz, \emph{Sparsifying the operators of fast matrix multiplication
  algorithms}, arXiv:2008.03759, 2020.
\bibitem{karstadt2017} E.~Karstadt and O.~Schwartz, \emph{Matrix multiplication,
  a little faster}, SPAA 2017; Journal of the ACM \textbf{67} (2020), no.~1.
\bibitem{holtz2025} O.~Holtz, J.~Hsu, S.~Moran, O.~Schwartz and N.~Wiernik,
  \emph{Alternative bases for new fast matrix multiplication algorithms},
  ACDA 2025.
\bibitem{karunaratne2026} S.~Karunaratne and A.~Idamekorala, \emph{55 additions
  suffice for $3\times 3$ matrix multiplication at rank 23},
  arXiv:2607.28676, 2026.
\bibitem{probert1976} R.~L.~Probert, \emph{On the additive complexity of matrix
  multiplication}, SIAM Journal on Computing \textbf{5} (1976), 187--203.
\bibitem{bshouty1995} N.~H.~Bshouty, \emph{On the additive complexity of
  $2\times 2$ matrix multiplication}, Information Processing Letters
  \textbf{56} (1995), 329--335.
\bibitem{plinopt} J.-G.~Dumas, B.~Grenet, C.~Pernet and A.~Sedoglavic,
  \emph{PLinOpt: C++ routines for linear, bilinear and trilinear straight-line
  programs}, \texttt{github.com/jgdumas/plinopt}.
\bibitem{mccormick1983} S.~T.~McCormick, \emph{A combinatorial approach to some
  sparse matrix problems}, Ph.D. thesis, Stanford University, 1983.
\bibitem{colemanpothen1986} T.~F.~Coleman and A.~Pothen, \emph{The null space
  problem I. Complexity}, SIAM Journal on Algebraic and Discrete Methods
  \textbf{7} (1986), 527--537.
\bibitem{tillmann2014} A.~M.~Tillmann and M.~E.~Pfetsch, \emph{The computational
  complexity of the restricted isometry property, the nullspace property, and
  related concepts in compressed sensing}, IEEE Transactions on Information
  Theory \textbf{60} (2014), 1248--1259.
\bibitem{sparsevectorfocs2025} \emph{Inapproximability of finding sparse vectors
  in codes, subspaces, and lattices}, FOCS 2025; arXiv:2410.02636.
\bibitem{qu2020} Q.~Qu, Z.~Zhu, X.~Li, M.~C.~Tsakiris, J.~Wright and R.~Vidal,
  \emph{Finding the sparsest vectors in a subspace: theory, algorithms, and
  applications}, arXiv:2001.06970, 2020.
\bibitem{vardy1997} A.~Vardy, \emph{The intractability of computing the minimum
  distance of a code}, IEEE Transactions on Information Theory \textbf{43}
  (1997), 1757--1766.
\bibitem{berlekamp1978} E.~R.~Berlekamp, R.~J.~McEliece and
  H.~C.~A.~van~Tilborg, \emph{On the inherent intractability of certain coding
  problems}, IEEE Transactions on Information Theory \textbf{24} (1978),
  384--386.
\bibitem{downey1999} R.~G.~Downey, M.~R.~Fellows, A.~Vardy and G.~Whittle,
  \emph{The parametrized complexity of some fundamental problems in coding
  theory}, SIAM Journal on Computing \textbf{29} (1999), 545--570.
\bibitem{dumer2003} I.~Dumer, D.~Micciancio and M.~Sudan, \emph{Hardness of
  approximating the minimum distance of a linear code}, IEEE Transactions on
  Information Theory \textbf{49} (2003), 22--37.
\bibitem{bhattacharyya2021} A.~Bhattacharyya, \'E.~Bonnet, L.~Egri, S.~Ghoshal,
  Karthik~C.~S., B.~Lin, P.~Manurangsi and D.~Marx, \emph{Parameterized
  intractability of even set and shortest vector problem}, Journal of the ACM
  \textbf{68} (2021), no.~3, article 16.
\bibitem{bennett2023} H.~Bennett, M.~Cheraghchi, V.~Guruswami and J.~Ribeiro,
  \emph{Parameterized inapproximability of the minimum distance problem over all
  fields and the shortest vector problem in all $\ell_p$ norms}, STOC 2023.
\bibitem{stephensdavidowitz2019} N.~Stephens-Davidowitz and V.~Vaikuntanathan,
  \emph{SETH-hardness of coding problems}, FOCS 2019, 287--301.
\bibitem{sanjose2025} R.~San-Jos\'e, \emph{An algorithm for computing generalized
  Hamming weights and the Sage package GHWs}, ACM Transactions on Mathematical
  Software \textbf{51} (2025), no.~4.
\bibitem{zimmermann1996} K.-H.~Zimmermann, \emph{Integral Hecke modules, integral
  generalized Reed--Muller codes, and linear codes}, Technical Report 3-96,
  Technische Universit\"at Hamburg-Harburg, 1996. (With A.~Brouwer's earlier
  work, the Brouwer--Zimmermann algorithm.)
\bibitem{lisonek2016} P.~Lison\v{e}k and L.~Trummer, \emph{Algorithms for the
  minimum weight of linear codes}, Advances in Mathematics of Communications
  \textbf{10} (2016), 195--207.
\bibitem{mckay} B.~D.~McKay, \emph{Isomorph-free exhaustive generation},
  Journal of Algorithms 26 (1998), 306--324.
\bibitem{givaro} The CASYS team, \emph{Givaro: exact arithmetic over
  $\GF(p)$ and over the rationals}.
\end{thebibliography}

\end{document}
